\section*{Chapitre \ref{ch:dv:dividends-borrow-and-forwards} --- Dividendes, emprunt et forwards}

\begin{solution}{exo:dv:dividends-borrow-and-forwards:1}
$(100-e^{-0.0075}-e^{-0.0225})e^{0.03}=101.0152$.
\end{solution}

\begin{solution}{exo:dv:dividends-borrow-and-forwards:2}
Pente $(-2.83-7.02)/10=-0.985$~: le facteur d'actualisation vaut 0.985. Puis $7.02=0.985(F-95)$ donne
$F=102.1269$.
\end{solution}

\begin{solution}{exo:dv:dividends-borrow-and-forwards:3}
$\bar D=4e^{-0.015}=3.9404$, $S/(S-\bar D)=1.0410$~;
$\sigma_{\mathrm{eff}}=0.25\sqrt{1.0410^2\times0.5+0.5}=0.2552$.
\end{solution}

\begin{solution}{exo:dv:dividends-borrow-and-forwards:4}
Chacune des deux options peut être exercée par anticipation, donc $C-P$ est seulement encadré~:
$S-D-K\le C-P\le S-Ke^{-rT}$ pour des options américaines sur un titre versant un dividende. Avant la date de détachement
le call porte une prime d'exercice anticipé (il peut être exercé pour capter le
dividende), si bien que $C-P$ est trop élevé et que le forward déduit de la relation européenne est trop
élevé~: le \omterm{def:dv:dividends-borrow-and-forwards:implied}{dividende implicite} ressort trop faible.
\end{solution}

\begin{solution}{exo:dv:dividends-borrow-and-forwards:5}
Dividende en numéraire~: $F=(S-4e^{-0.015})e^{0.03}$, donc $\partial F/\partial S=e^{0.03}=1.0305$.
\omterm{def:dv:dividends-borrow-and-forwards:prop}{Dividende proportionnel}~: $F=0.96\,Se^{0.03}$, donc $\partial F/\partial S=0.9892$. Le même dividende
donne deux couvertures distantes de 4~\%.
\end{solution}

\begin{solution}{exo:dv:dividends-borrow-and-forwards:6}
Le forward monte de 100.3611 à 101.1038 et le call de 12.7558 à 13.1475~: $+0.3917$,
soit environ 3~\% de sa valeur, pour une baisse de 10~\% des dividendes.
\end{solution}

\begin{solution}{exo:dv:dividends-borrow-and-forwards:7}
Forward 100.1132, facteur d'actualisation 0.97033~; avec les dividendes connus (valeur actuelle 2.3661)
le \omterm{def:dv:dividends-borrow-and-forwards:implied}{taux d'emprunt implicite} est de 0.504~\%, contre 0.5~\%~: l'erreur de régression, d'environ un centime sur le
forward, représente 0.004~\% du taux d'emprunt.
\end{solution}

\begin{solution}{exo:dv:dividends-borrow-and-forwards:8}
Les deux volatilités sont calculées sur le mauvais forward. Sur le forward implicite dans les prix des options,
48.74, le call et le put à 50 impliquent tous deux 60~\%~: il n'existe aucun écart de volatilité à
capter. La stratégie est un renversement déguisé~: vente de puts et achat de calls, couverte par
une vente à découvert d'actions, qui nécessite un emprunt à 35~\%, ce qui coûte ce que le «~edge~» semble rapporter.
\end{solution}

\begin{solution}{pb:dv:dividends-borrow-and-forwards:1}
\textbf{1.} 3.73, 1.23, $-1.26$, $-3.75$, $-6.24$.
\textbf{2.} Facteur d'actualisation 0.99680~; forward 48.74.
\textbf{3.} 3.90~\%~; les cotations sont arrondies au centime, et une erreur d'un centime sur une pente
mesurée sur dix points de strike déplace le facteur d'actualisation de $10^{-3}$.
\textbf{4.} 1.42.
\textbf{5.} 35.0~\% par an.
\textbf{6.} 50.16.
\textbf{7.} Call 3.50, put 3.34.
\textbf{8.} Le call varie de $+0.72$ (26~\% de son prix), le put de $-0.70$.
\textbf{9.} Call 47.3~\%, put 72.3~\%.
\textbf{10.} La parité relie $C-P$ au forward~; avec le forward naïf, les volatilités du call et du
put doivent différer pour l'ajuster. Sur le forward 48.74, les deux valent 60~\%.
\textbf{11.} $51.26-50e^{-0.04\times30/365}$~: 1.42 aujourd'hui.
\textbf{12.} $50\times0.35\times30/365=1.44$.
\textbf{13.} Non~: le profit apparent est le coût de l'emprunt, à l'arrondi des cotations près.
\textbf{14.} Vendre l'action et la racheter synthétiquement (acheter le call, vendre le put)~:
cela rapporte 1.42 sur le mois, comme le prêt de l'action à 35~\%. Les détenteurs qui
ne peuvent pas prêter (un fonds dont le dépositaire ne prête pas) perçoivent ainsi ces frais.
\textbf{15.} Acheter des puts ou des put spreads, ou vendre des calls~: les options portent le coût de l'emprunt dans
leurs prix, si bien que le pari est coûteux mais exprimable.
\textbf{16.} Le forward bondit de 48.74 à 49.96~: le call 50 monte et le put baisse
d'environ 0.6 chacun, un mouvement sans lien avec la volatilité ni le spot.
\textbf{17.} Les frais d'emprunt agissent comme un rendement de dividende~: un call américain peut avoir intérêt à
l'exercer par anticipation pour détenir l'action et la prêter au taux des frais, ce qui ajoute une
prime d'exercice anticipé aux calls et brouille davantage la parité.
\textbf{18.} Le forward implicite dans les options elles-mêmes, recalé sur le marché de l'emprunt,
et non le spot multiplié par le taux de financement.
\textbf{19.} 35.0~\% par an~; le prix naïf du call 50 est trop élevé de 0.72.
\textbf{20.} Le forward et, à travers lui, le coût de la vente à découvert de l'action.
\end{solution}

\begin{solution}{iq:dv:dividends-borrow-and-forwards:1}
Parité call-put~: $C-P=P(0,T)(F-K)$. Régresser $C-P$ sur $K$ à partir des cotations européennes proches de la
monnaie~; la pente vaut $-P(0,T)$, l'ordonnée à l'origine $P(0,T)F$. Avec des options américaines, utiliser
les cotations hors de la monnaie, où l'exercice anticipé est négligeable, ou retirer d'abord sa prime.
\iqlookfor{la parité comme régression linéaire sur l'ensemble des strikes.}
\end{solution}

\begin{solution}{iq:dv:dividends-borrow-and-forwards:2}
Le forward est inférieur au spot~: un dividende important avant l'échéance, ou des frais d'emprunt élevés (un
titre difficile à emprunter). Les deux entrent dans la parité comme du carry~; au forward, les deux options sont
de même valeur.
\iqlookfor{dividendes et emprunt à travers le forward.}
\end{solution}

\begin{solution}{iq:dv:dividends-borrow-and-forwards:3}
En numéraire pour les dividendes proches, déjà annoncés~; proportionnel (ou un mélange) pour les lointains,
car le dividende d'une société baisse quand son action baisse. La couverture change parce que la
sensibilité du forward au spot vaut $1/P(0,T)$ avec des dividendes en numéraire et $F/S$ avec
des \omterm{def:dv:dividends-borrow-and-forwards:prop}{dividendes proportionnels}~; pour une option à cinq ans, l'écart est de l'ordre du rendement de dividende
multiplié par cinq ans.
\iqlookfor{la dynamique, pas seulement le niveau du forward.}
\end{solution}

\begin{solution}{iq:dv:dividends-borrow-and-forwards:4}
La sensibilité des positions aux dividendes attendus. Les émetteurs de produits indexés sur actions
structurés, dont les couvertures perçoivent les dividendes promis aux détenteurs de notes,
sont acheteurs~; ils vendent des \omterm{def:dv:dividends-borrow-and-forwards:risk}{swaps de dividendes} et des contrats à terme. En 2020, les superviseurs ont demandé aux banques de suspendre
les versements de dividendes, tandis que les dividendes d'indice se négociaient bien en deçà des modèles et que les ventes de couverture des émetteurs
ont contribué à la chute.
\iqlookfor{qui le porte, et pourquoi la chute de 2020 a été amplifiée.}
\end{solution}

\begin{solution}{iq:dv:dividends-borrow-and-forwards:5}
Options~: \omterm{def:dv:dividends-borrow-and-forwards:escrowed}{modèle à dividendes séquestrés} (un arbre sur $S-\mathrm{PV}(D)$, en rajoutant la valeur actuelle pour l'exercice), ou une
baisse des valeurs aux nœuds de $D$ à la date de détachement. L'arbre séquestré sous-estime la
volatilité avant la date de détachement, à moins que sa volatilité ne soit ajustée~; une baisse de $D$ casse
la recombinaison (l'arbre doit être reconstruit après la date de détachement ou interpolé), et une
baisse proportionnelle conserve la recombinaison mais est fausse pour un dividende en numéraire. Tester contre une
référence européenne de même forward.
\iqlookfor{le problème de recombinaison et le décalage de volatilité.}
\end{solution}

\begin{solution}{iq:dv:dividends-borrow-and-forwards:6}
Le modèle sur spot~: sa volatilité s'applique à l'ensemble de l'action, y compris la part versée,
de sorte qu'avant la date de détachement, le forward bouge davantage. Pour une option d'un an avec un dividende de 4~\% versé à
six mois, environ un demi-point de volatilité, 0.2 de prime à la monnaie~; l'ajustement est
$\sigma^2(S/(S-\bar D))^2t_1+\sigma^2(T-t_1)$.
\iqlookfor{de quoi l'on mesure la volatilité.}
\end{solution}
